\section{The Postconciliar Reform: Mandate, Construction, Promulgation}

The postconciliar Missal is neither a conciliar committee's private experiment nor a text composed by the fathers of Vatican II. It has a layered authorship. \emph{Sacrosanctum Concilium} supplied binding principles and concrete mandates. Paul VI created the Consilium to implement them. Expert groups researched, drafted, tested, and revised. The competent congregations issued instructions. The pope reserved and decided major questions, promulgated the Order and Missal, and required their use. Episcopal conferences translated and adapted within law. Parish practice then added a final, often uncontrolled, layer that must never be confused with the typical book.

\subsection{The council's positive commands and limiting grammar}

The constitution calls the liturgy the exercise of Christ's priestly office and the summit and source of ecclesial life. Reform serves that theological reality. It asks for rites of noble simplicity whose parts and connections can be perceived; a richer table of Scripture; restored common intercession; homily and liturgical instruction; extension of vernacular within competent authority; widened concelebration and Communion under both kinds in stated cases; and the full, conscious, active participation of the baptized.

Its restraints are equally textual. Regulation belongs to the Apostolic See and, within law, bishops and conferences; no priest may change the rite on his own. Latin is to be preserved in the Latin rites. Gregorian chant has principal place. Before innovation there must be careful theological, historical, and pastoral investigation; there must be a genuine and certain need; and new forms are to grow organically from existing forms (SC 22--23, 36, 50--58, 116). These are not traditionalist additions to the council. They are the council's criteria for judging implementation.

\begin{fourstable}{Question}{Council actually said}{Later lawful development}{Frequent false attribution}
Language & Latin is preserved; vernacular may be extended by competent authority. & Conferences and Rome authorized almost complete vernacular books in many territories. & “Vatican II abolished Latin.”\\
Orientation & The constitution gives no command to turn the celebrant toward the people throughout Mass. & Freestanding altars and versus-populum celebration spread through later instructions, design, and custom; the current Missal does not make that posture an element of validity. & “The council ordered every altar turned around.”\\
Communion in the hand & The constitution is silent. & \emph{Memoriale Domini} and later territorial recognitions created the present regulated permission. & “The council restored Communion in the hand.”\\
Music & Gregorian chant has principal place, while other sacred music is admitted. & Conferences, composers, and parishes developed a wide vernacular repertoire. & “The council abolished chant and required contemporary popular idioms.”\\
\end{fourstable}

\subsection{Consilium and papal responsibility}

Paul VI established the Consilium in 1964 under Cardinal Giacomo Lercaro, with Annibale Bugnini as secretary. Study groups addressed the Order of Mass, lectionary, collects, prefaces, Eucharistic prayers, calendar, chant, and the other sacramental books. Their sources included ancient Roman sacramentaries, medieval and later Roman texts, non-Roman Western and Eastern witnesses, Scripture, patristic scholarship, pastoral experiments, and new composition.

Calling the result “fabricated” hides the received materials; calling it merely an organic pruning hides the method and scale. The Roman Canon survived as Eucharistic Prayer I. The Prayer of the Faithful and a broader Old Testament lectionary recover ancient patterns. Many collects and prefaces have venerable sources. Yet prayers were edited, combined, relocated, or newly written; an entirely new multi-year lectionary was built; three principal Eucharistic prayers joined the Canon; the offertory, calendar, rubrics, and ministerial ecology were extensively reconstructed. It is accurately a papally promulgated, committee-developed, source-rich twentieth-century synthesis.

Bugnini's memoir is indispensable participant testimony, not an official transcript immune from criticism. Allegations of a proved Masonic authorship do not follow from the public record and are not needed for a strong historical critique. Paul VI reviewed reports, settled disputed points, spoke publicly about the sacrifice demanded by change, and promulgated the result. Ultimate juridical responsibility is papal, not a device by which every disputed choice can be assigned to unnamed experts.

\subsection{The reform arrived in stages}

\begin{longtable}{>{\bfseries\raggedright\arraybackslash}p{0.15\linewidth}>{\raggedright\arraybackslash}p{0.28\linewidth}>{\raggedright\arraybackslash}p{0.49\linewidth}}
\toprule
\textbf{Date} & \textbf{Stage} & \textbf{What changed or became authoritative}\\
\midrule
\endhead
1964 / 1965 & \emph{Inter Oecumenici} and transitional books & Simplifications and wide vernacular use altered the enacted 1962 order before a complete replacement Missal existed. The “1965 Mass” is a transitional state, not the 1970 typical edition.\\
1967 & \emph{Tres abhinc annos}; experimental “normative Mass” & Further changes and a synodal demonstration informed continuing work. Neither event by itself promulgated the final Order.\\
1968 & Three additional Eucharistic prayers & Ended the Roman Canon's practical exclusivity while retaining it as the first prayer.\\
3 Apr 1969 & Paul VI, \emph{Missale Romanum} & Promulgated the restored Missal, naming the new Order, expanded readings, prayers, and calendar, with implementation beginning at Advent.\\
1969 & \emph{Ordo Missae}, calendar, and lectionary & Supplied the new ritual spine and reading architecture while the complete book was still in production.\\
1970 & First Latin typical edition & Published the complete normative Missal. Doctrinal clarification of the General Instruction's opening description answered criticism of its first wording.\\
1975 & Second Latin typical edition & Corrected and developed the postconciliar book; it was not a new rite unrelated to 1970.\\
2002 / 2008 & Third Latin typical edition and emended reprint & Added saints, formularies, rubrics, and other revisions. These are the Latin witnesses used for the present comparison.\\
2011, United States & Current English implementation & Replaced the 1973-era translation with a more formally corresponding rendering of the third typical edition. Territorial texts and adaptations remain subordinate to the Latin edition and approved law.\\
\bottomrule
\end{longtable}

Promulgation, publication, translation, and reception were therefore not one event. The constitution took legal effect before every complete vernacular altar book existed. National conferences worked at different speeds. Provisional translations overlapped with new rubrics. Architecture, music, catechesis, experimentation, and the cultural revolution of the late 1960s made the experienced reform broader than the printed Missal.

\subsection{Replacement, abrogation, and the older book}

The revised Missal was distributed as the normative Roman book, not as a boutique option. A 1974 notification stated that, apart from specified concessions, only the revised Missal could be used once approved vernacular editions appeared. Limited permissions followed, including a 1971 concession in England and Wales for certain celebrations, the 1984 indult \emph{Quattuor abhinc annos}, and the widened pastoral provisions after \emph{Ecclesia Dei} in 1988.

Benedict XVI judged in \emph{Summorum Pontificum} (2007) that the 1962 Missal had never been juridically abrogated and regulated it as an “extraordinary form” alongside the ordinary form. Francis's \emph{Traditionis custodes} (2021) changed that juridical framework, declared the postconciliar books the unique expression of the Roman Rite's lex orandi, and abrogated inconsistent prior norms. The historical fact of practical replacement, Benedict's legal judgment about formal abrogation, and Francis's present restrictive law are three different propositions.

\subsection{Continuity, gain, and cost}

The new book preserves the essential structure of gathering, Word, Eucharistic offering, Communion, and mission. It confesses sacrifice, Real Presence, ministerial priesthood, intercession for living and dead, and the communion of saints. Its lectionary proclaims vastly more Scripture; its prefaces and formularies enlarge the Church's audible theology; its responses and intercessions can reveal the assembly's differentiated participation.

The cost is not imaginary. A dense inherited ecology was changed simultaneously: the one-year cycle, threshold prayers, medieval offertory, silent and exclusive Canon, repeated private Communion preparation, stable calendar memory, and sharply limited presidential choice. The practical predominance of vernacular, versus-populum celebration, amplified speech, modern music, and lay sanctuary ministries intensified the perceptual break even where not commanded by the council. Rapid implementation and unauthorized experiment made it harder to distinguish reform from rupture.

\governingjudgment{Vatican II genuinely required reform; it did not write the Missal or command every practice later associated with it. Paul VI validly and lawfully promulgated a materially reconstructed Roman book. Catholic fidelity affirms that authority. Historical fidelity still asks, under the council's own organic-growth criterion, which changes were necessary, which were fruitful, and which carried avoidable losses.}
